% !TeX root = main.tex
% !TeX program = xelatex
% Definitions are tied to the archived inputs identified in the source comments.

\section{形式化建模与验证设置}
\label{sec:formal-modeling-analysis-setup}

规则选择、带身份的凭证消费和共享额度更新分别决定 SAGA 授权约束的执行。
相应分析以可执行参考语义、带事件的多重集重写规则和可交错状态转换，
刻画预算决定、凭证使用身份与累计请求次数。

\subsection{建模方法与抽象原则}
\label{subsec:modeling-approach-abstractions}

协议规定确定预期授权语义，参考实现确定状态更新与并发操作的执行结构。
分析依据固定版本的 SAGA 代码~\cite{sagaReferenceImpl2026}，
提取匹配器的规则选择、Provider 的按身份发放状态、接收方的密钥消费及令牌检查与扣减。
确定性的预算求值通过独立参考计算与生产函数对照，
多方 OTK 交互采用 Tamarin 的多重集重写表示~\cite{meier2013tamarin}，
共享 ACT 额度采用 TLA$^{+}$ 状态与动作表示~\cite{lamport2002specifying}。
这些表示分别保留对应性质所依赖的决策、身份或交错行为。

源码对应的依据为人工映射与选定的受控执行，OTK 符号分析结果在理想原语和信道假设下解释。
形式化实现精化和计算安全转移尚未建立。
模型、完整查询、源码映射和原始运行记录按固定工件版本提供\footnote{\raggedright
\href{https://github.com/g1020268262-a11y/saga/tree/13455bb644f954964ef2fcd9303f50d0fdc43830}{\texttt{github.com/g1020268262-a11y/saga}},
revision \texttt{13455bb6}.
OTK 和 ACT 工件分别位于 \texttt{proofs/otk\_transfer} 与 \texttt{proofs/quota\_validation}。}。

\subsection{策略求值模型}
\label{subsec:policy-evaluation-model}

参考策略语义以匹配规则的具体性确定预算。
对接触策略 $CP_A$ 与发送方身份 $s$，参考端独立计算匹配与评分，并返回唯一最高评分匹配规则的预算。
支持的语法为精确小写 ASCII 身份、字面 UID 通配及全通配。
生产端直接调用原始 \texttt{match()}，检查目标为
\begin{equation}
  \operatorname{Eval}_{\mathrm{impl}}(CP_A,s)
  =\operatorname{Eval}_{\mathrm{ref}}(CP_A,s).
  \label{eq:policy-evaluation-consistency}
\end{equation}
两侧分别表示实际函数与独立参考求值的预算。
对照采用相同身份、相同规则及不同规则排列。
Provider 的 \texttt{/access} 路由与接收方的
\texttt{handle\_i\_agent\_\allowbreak connection()} 均调用生产匹配器，
消费者执行进一步记录两处决策及其后的凭证状态。

% Reference evaluator: proofs/decision_bridge/bridge.py:49--93.
% Consumer: proofs/consumer_validation/harness.py:102--113.

\subsection{OTK 发放与使用模型}
\label{subsec:otk-issuance-use-model}

OTK 的发放、持有与消费分别对应具有不同更新位置的资源状态。
模型固定 Provider $P$、接收方 $A$ 和身份独立的发送方 $B,C$，
设置两枚 OTK，以及每个发送方一次不可补充的领取额度。
线性事实 \texttt{ProviderCredit} 表示领取机会，
\texttt{Provider\allowbreak UnusedOTK} 与 \texttt{Receiver\allowbreak UnusedOTK}
分别表示尚未发放和尚未消费的密钥。
持久事实 \texttt{HeldOTK} 表示已持有的公钥材料。
Provider 的单次数据库发放更新对应同时消费领取机会与待发放 OTK 的规则，
接收方本地密钥表的检查与删除则对应另一条消费规则。
因此，发放改变 Provider 的资源状态，后续 ACT 签发才改变接收方的密钥状态。

% Source map: proofs/otk_transfer/SOURCE_MAP.zh-CN.md.
% Rules: proofs/otk_transfer/saga_otk_transfer.spthy.
协议交互通过保留实际连接者身份的规则继续执行。
已登记身份与 PAC 的绑定由初始持久事实表示，TLS 以带端点身份的私有消息事实表示。
接收方核对本地许可和未消费 OTK，使用其私钥与连接者 PAC 派生密钥，并签发绑定该 PAC 的 ACT。
发送方保留对应的 DH 密钥依赖，正确解密后才能提交模型中的应用请求。
DH 使用内建等式理论，KDF 以带域标签的哈希项表示，
ACT 以包含新鲜标识及 PAC 的理想对称密文表示。

主动转交规则允许 $B$ 将已领取的 OTK 公钥发送给 $C$，
后续签发仍受接收方的未消费状态约束。
$C$ 使用自己的身份、公钥和私钥继续协议。
主配置中，两发送方均获 $P$ 与 $A$ 的策略允许，长期私钥保持保密。
\texttt{UniqueSetup} 约束单次初始化，每枚已领取 OTK 至多触发一次转交，访问与令牌申请规则可以重复触发。
运行采用 Tamarin 1.12.0 与 Maude 3.5.1，
检查允许转交、关闭转交及两端拒绝 $C$ 三种配置。
原始轨迹通过事件投影核对发放、解密和准入的先后关系。
独立证明证书检查尚未实施。

% Archive: proofs/otk_transfer/evidence/run-20261025T033815645931Z/.
发放数量与使用身份分别由计数性质和事件对应刻画。
记 $I_{A,s}$ 为执行前缀中 $P$ 向获准身份 $s$ 发放 $A$ 的 OTK 的次数，
$\beta_{A,s}$ 为其预算，原有发放约束为 $I_{A,s}\leq\beta_{A,s}$。
$\mathsf{Issue}(A,s,k)$ 表示向 $s$ 发放 $k$ 的事件。
$\mathsf{Use}(A,s,k)$ 表示匹配的新 ACT 已签发、由 $s$ 正确解密且其首个应用请求获准入的事件投影。
对任意 $s,k$ 和位置 $i$，额外检查
\begin{equation}
  \begin{aligned}
    &\mathsf{Use}(A,s,k)@i\\
    &\quad\Rightarrow\exists j<i.\ \mathsf{Issue}(A,s,k)@j.
  \end{aligned}
  \label{eq:otk-same-sender-correspondence}
\end{equation}
该对应关系是将领取预算解释为按身份限制新 ACT 取得次数时的更强目标。
每枚 ACT 的使用观察到其首次请求准入。

\subsection{ACT 并发配额模型}
\label{subsec:act-quota-model}

共享额度的并发表示由源码中两段临界区的边界确定。
模型从同一有效 ACT 已在双端保存后开始，
保留双端余额、连接局部位置、消息交付状态及准入与后端调用计数。
状态变量 \texttt{receiverRemaining} 记录接收方余额 $r$，
观测算子 \texttt{Accepted} 对各连接的准入计数求和。
\texttt{ACheck} 对应 \texttt{token\_is\_valid()} 的检查临界区，
读取余额并使连接进入已检查位置，随后释放锁。
\texttt{ADebit} 对应 \texttt{receive\_conversation()} 中重新加锁的扣减，
更新余额为 $\max(0,r-1)$、增加准入计数并推进连接位置。
扣减动作以连接已经完成检查为前提。
后续后端处理由 \texttt{ABackend} 记录。
该划分保留连接在检查与扣减之间的调度位置~\cite{sagaReferenceImpl2026}。

% Model: proofs/quota_validation/SagaQuotaDelivery.tla:132--160,290.
% Config: proofs/quota_validation/evidence/delivery-20261024T150331844425Z/inputs/delivery_configs/concurrent_race.cfg.
累计准入上限以进入处理流程的请求数量为判断依据。
对初始额度为 $Q(t)$ 的令牌 $t$，记累计准入为 $N_{\mathrm{adm}}(t)$，检查
\begin{equation}
  N_{\mathrm{adm}}(t)\leq Q(t).
  \label{eq:act-admission-bound}
\end{equation}
$N_{\mathrm{adm}}(t)$ 对应 \texttt{Accepted}，余额表示扣减后的可用状态，后端调用另行计数。
TLC 2.19 在初额 $1$、两个连接各至多一个普通请求、
非配额元数据有效且零故障的配置中检查可达状态不变量。
精确输入、工具版本、日志及输出状态随运行记录归档。
